% ============================================================
\chapter{Results}
\label{chap:results}
% ============================================================

The evaluation asks six questions. Does content-addressed caching make repeated
submissions cheap (Section~\ref{sec:res-cache})? What does interception cost
when it decides not to distribute (Section~\ref{sec:res-shim})? Does parallel
work cross machines, and is it shared well across unequal ones
(Section~\ref{sec:res-dist})? Does training across machines stay accurate, and
when does it pay (Section~\ref{sec:res-train})? Does a run survive the loss of a
machine (Section~\ref{sec:res-ft})? And what happens between real machines on
different networks (Section~\ref{sec:res-recorded})?

\section{Evaluation Setup}
\label{sec:res-setup}

\textbf{Hosts.} Laboratory measurements ran on two development hosts: an Intel
i5-13420H with 12 hardware threads and 15\,GiB of memory under Fedora~43, and a
16-thread laptop with an RTX~3060 laptop GPU under Fedora~44.

\textbf{Container labs.} Worker daemons run in containers on one host.
In the \emph{plain lab} (three workers) the containers share the host's
processor, memory and loopback interface; it measures correctness, overhead,
cache behaviour and fault handling, and cannot measure multi-machine speedup.
Where a figure depends on speed or on the network, each container is given a
fixed processor and memory limit and the links between containers are shaped by
\texttt{tc netem} to the latency, bandwidth and loss of a LAN, a Wi-Fi network
or a 5G link, with up to 5\% loss; speedup is then one container against $N$
containers of the same size. The \emph{LAN lab} reproduces a college
laboratory: containers on one bridge network, the host's own daemon as the GPU
machine and one container on multi-user Nix, with a two-container variant
sharing the host's GPU for training rings.

\textbf{Two-site testbed.} Joining across networks without Tailscale is
evaluated on a two-site container testbed on the same host: two containers on
separate private networks, both numbered 192.168.0.0/24, each behind a router
container whose iptables rules translate addresses and drop unsolicited inbound
packets, with the introducer on a third network standing for the internet. Every
user command is run across it.

\textbf{Training accuracy.} The accuracy of each synchronisation setting is
measured on three models --- a four-layer perceptron sized for a GPU (MLP-G), the
same perceptron at a CPU sizing (MLP-C), and a small transformer --- with three
seeds per cell, against one device trained on the same data: two-rank rings in
the lab, larger rings in-process, with the ranks side by side on a single machine.

\textbf{Real machines.} Two laptops in different homes (RTX~3060 and RTX~3050)
connected over a direct Tailscale path \cite{anderson2020nat}, a Samsung
Galaxy S20~FE phone under Termux, and further measurements between a 20-core
machine and a 16-core peer joined over Tailscale, with a public cloud host
acting as introducer.

\textbf{Workloads.} Seven example workloads cover a standard-library parallel
map, a map over a lambda, scikit-learn with \texttt{joblib}, NumPy linear algebra,
out-of-core pandas, torch data-parallel training and file synchronisation; one
script combines every kind of intercepted work, and five idiomatic, unoptimised
user scripts stand in for unseen workloads. Lab
wall times are medians of three to five runs, and accuracy is the mean of three
seeds.

\section{Functional Verification}
\label{sec:res-tests}

The test suite exercises the properties the design relies on, with every
optional Python library present --- among them that a hash-shuffled aggregation
equals local pandas, that a parallel map whose every remote part fails still
returns the correct result, that the daemon's average is weighted by samples,
and that a link hello replayed on another session is refused. A mutation audit
guards against a suite that passes without proving anything: it breaks each of
95 features in turn, from the sample-weighted gradient fold to the
router-mapping checks, and requires a test to fail. Every mutation was caught.

\section{Environment Cache}
\label{sec:res-cache}

The same trivial job was submitted five times to the plain lab
(Figure~\ref{fig:rescache}). The first submission, which also built the
environment, took 23.4\,s; the second landed on a worker that did not yet hold
the closure and took 8.1\,s; once the closure was in place, a submission sent
only the workspace and took a median of 0.83\,s.

\begin{figure}[!htbp]
    \centering
    \includegraphics[width=87mm]{\figpath{fig_5.1_closure_cache.png}}
    \caption[Cold and warm submissions]{Left: first submission against the median of the submissions that
    sent nothing. Right: every run, coloured by what its log says was sent.}
    \label{fig:rescache}
\end{figure}

Between two real machines, the path-level difference sent 29\,MB instead of
206\,MB when one package was added to an environment the peer held, and
0.43\,MiB of a 444\,MiB closure for an environment that differed from the held
one by one small pure-Python package. An unseen script's environment built in 4\,s on a machine that
had built the common packages ahead of time, against 54\,s on one that had not.

With one of two nodes held to 800\,kbit/s, the first run waited the upload's
60\,s limit for the slow node and then started; the background send to that
node took about 15 minutes and held up no run, and once it had finished a run
took 2.1\,s with the slow node taking about a third of the items. With the
cache cap lowered, both caches were pruned and the next job still ran on its
peer.

\section{The Cost of Leaving Interception On}
\label{sec:res-shim}

A 16-item parallel map whose tail is 0.01\,s of work against a 0.5\,s dispatch
floor is declined by the cost model; a forced run of the same workload does
dispatch, so the decline is the model's choice and not a dead cluster.
Declining costs about 0.12\,s per job, the measurement of the head of the map
plus the decision: a median of 0.80\,s against 0.68\,s with interception
disabled. Interpreter start-up for a job that imports nothing heavy is
$1.24\times$ that of plain Python when NumPy, pandas, torch and joblib are
patched on first import, against $20.0\times$ when they are patched at
start-up (Figure~\ref{fig:resshim}). A continuous-integration check fails any
change that raises this ratio above $3\times$.

\begin{figure}[!htbp]
    \centering
    \includegraphics[width=150mm]{\figpath{fig_5.2_shim_overhead.png}}
    \caption[Cost of interception]{Left: the declined 16-item parallel map, median
    and range; the enabled side's cold first run is left out of its range. Right: interpreter
    start-up with libraries patched at start-up and on first import.}
    \label{fig:resshim}
\end{figure}

\section{Distribution Across Machines}
\label{sec:res-dist}

\textbf{Work crosses nodes.} With distribution forced, over a third of the
items of a 48-item parallel map ran on other daemons of the plain lab, with
results identical to a local run. The NumPy workload moved its 34\,MB singular value decomposition to one
worker. On one host distribution can only add overhead, and the pool benchmark shows
how much: plain Python took 0.17\,s, a hand-tuned pool 0.17\,s, and the
untouched script under PipedPeer 0.24\,s with a quarter of the items on the
cluster.

\textbf{Sharing by measured speed.} On a lab cluster limited to 8, 2 and
1 CPUs, dividing a map by headcount ran 1.6 times slower than not distributing at
all. Dividing by measured throughput declined the one-CPU worker and took
2.9\,s, against 3.6\,s by headcount. With one peer held to a tenth of a
processor in a cluster of otherwise equal peers, the slowdown against an equal
cluster fell from 1.6 times on the first run to 1.08 by the fourth, as the fit
learned the slow peer.

\textbf{Scaling.} A scaling map sweeps one variable at a time with the rest
fixed --- payload from 1\,KB to 64\,MB, 8 to 2048 items, per-item cost from
$10^3$ to $3\times10^6$ operations, two to four nodes, and a speed spread of up to
ten times --- and fits each step's growth on a log-log scale from the run
records. No step grew faster than its expected order, and four equal nodes ran a
map 2.9 times faster than one (0.81\,s against 2.37\,s).

\section{Training}
\label{sec:res-train}

\textbf{One GPU or a split.} The decision between a ring and one GPU matters in
both directions. In the lab, a small model with heavy steps on a LAN
trained in 9.1\,s split across two ranks against 11.8\,s on one GPU, while a
9\,MB model with light steps on Wi-Fi took 80.7\,s split against 8.0\,s on one
GPU. On a single machine, the same model and batch took 1.5\,s on its GPU and 55.0\,s
on its CPUs.

\textbf{Fast synchronisation.} Table~\ref{tab:fastsync} and
Figure~\ref{fig:ressync} compare a two-rank ring under exact and fast
synchronisation every 64 steps with one device. Every model stays within the
1\% bar on both link profiles, each gap within the spread of exact
synchronisation's own gap, while each rank sends 9 to 16 times fewer bytes. Fast synchronisation was
faster than exact on every model and profile, by 1.2--1.4 times on the LAN and
1.6--3.1 times on Wi-Fi.

\begin{table}[!htbp]
\centering
\caption{Two-rank ring against one device, fast synchronisation every 64
steps: validation-loss gap (mean of three seeds), bytes sent per rank, and wall
time (one seed)}
\label{tab:fastsync}
\renewcommand{\arraystretch}{1.2}
\rowcolors{2}{}{rowgray}
\small
\begin{tabularx}{\textwidth}{X C{2.3cm} C{2.3cm} C{2.6cm} C{2.6cm}}
\toprule
\textbf{Model} & \textbf{Gap, fast LAN / Wi-Fi} & \textbf{MiB per rank, exact / fast} & \textbf{Seconds, exact LAN / Wi-Fi} & \textbf{Seconds, fast LAN / Wi-Fi} \\
\midrule
MLP-G (one device 27.9\,s) & $+0.64\%$ / $+0.61\%$ & 81.1 / 9.1 & 20.0 / 41.8 & 15.2 / 18.9 \\
MLP-C (one device 16.9\,s) & $-2.14\%$ / $-2.09\%$ & 98.9 / 10.7 & 14.6 / 43.7 & 10.4 / 14.1 \\
Transformer (one device 26.7\,s) & $+0.18\%$ / $+0.50\%$ & 21.7 / 1.4 & 22.7 / 35.9 & 19.1 / 22.9 \\
\bottomrule
\end{tabularx}
\end{table}

\begin{figure}[!htbp]
    \centering
    \includegraphics[width=150mm]{\figpath{fig_5.4_sync_modes.png}}
    \caption[Exact and fast synchronisation]{Left: wall time of one device and of a two-rank ring under exact and
    fast synchronisation on Wi-Fi. Right: validation-loss gap to one device under
    fast synchronisation; the dashed line is the 1\% bar.}
    \label{fig:ressync}
\end{figure}

\textbf{The outer step and the ring size.} DiLoCo's outer step (Nesterov momentum
0.9 at an outer learning rate of 0.7) failed the bar at every interval tried,
leaving the transformer $21.7\%$ above one device even with one synchronisation
per epoch (156 local steps), the longest interval these runs allow. In-process
runs on the same three models, on up to six seeds, set the interval $H(n)$ as
the largest that passed on every model: 32 steps at two ranks (64, which passed
in the lab above, failed the transformer in this harness on other seeds,
$+2.13\%$), 16 at three and 4 at four; at eight ranks the transformer failed at
every interval down to~2, so no interval is used for five or more ranks. Of ten methods tested at 4, 8 and 11 ranks, only exact synchronisation
with rank-16 PowerSGD \cite{vogels2019powersgd} compression and error feedback
kept every model within 1\% (worst mean $+0.18\%$), at $3.2\times$ fewer bytes
than 8-bit exact synchronisation on the perceptrons
(Appendix~\ref{app:extra}). Fast synchronisation therefore serves the two to
four machines of a small group.

\textbf{Adapting the ring.} On a ring of 16, 8 and 2 cores, admission took the 16
and the 8, shared the work 66:34, and ran in 23.6\,s against 25.4\,s with equal
batches. Throttling one rank from 8 CPUs to 1 and back moved the shares from
64:36 to 84:16 and then to 70:30.

\section{Fault Tolerance and Recovery}
\label{sec:res-ft}

Table~\ref{tab:drills} lists the failure drills run in the LAN lab and its
two-container GPU variant. The tests add that a parallel map whose every remote chunk fails
still returns the correct result (Section~\ref{sec:res-tests}).

\begin{table}[!htbp]
\centering
\caption{Failure drills}
\label{tab:drills}
\renewcommand{\arraystretch}{1.2}
\rowcolors{2}{}{rowgray}
\begin{tabularx}{\textwidth}{X C{5.2cm}}
\toprule
\textbf{Drill} & \textbf{Result} \\
\midrule
A node stopped during a parallel map & same checksum, \textbf{16.5}\,s against \textbf{17.0}\,s undisturbed \\
A worker killed during a remote job, two machines & node named gone \textbf{16}\,s after the kill \\
A training rank killed during a run & survivor finished all 80 epochs, \textbf{90.5}\,s \\
A training rank's machine back during a run & rejoined at the next epoch, \textbf{4}\,s after its daemon started \\
Interrupt during a parallel map & peers stopped within \textbf{2}\,s \\
Terminal closed during a run & stopped everywhere, as an interrupt \\
Detach, then attach while running and after it finished & followed live; replayed in \textbf{1.0}\,s once finished \\
Submitter's network removed during a remote job & \textbf{200}\,s outage: job ran once, results back \textbf{15}\,s after reconnect \\
Introducer's address changed, no commands run & recovered in about \textbf{12}\,s \\
Introducer suspended & live peer still reached in \textbf{0.03}\,s \\
\bottomrule
\end{tabularx}
\end{table}

\section{Sandbox Cost}
\label{sec:res-sandbox}

\begin{figure}[!htbp]
    \centering
    \includegraphics[width=101mm]{\figpath{fig_5.3_sandbox_overhead.png}}
    \caption[Sandbox start cost]{Sandbox start cost, median of 20 iterations.}
    \label{fig:ressandbox}
\end{figure}

The OCI runtime costs 15 to 20\,ms more per job than the lighter
\texttt{bwrap}, small beside a warm job of under a second, and its device model
is what makes GPU passthrough possible (Figure~\ref{fig:ressandbox}).

\section{End-to-End Runs in the LAN Lab}
\label{sec:res-lablayout}

The whole system was run end to end in the LAN lab for the processor steps and
its two-container GPU variant for the GPU steps. The containers
share one laptop, so Table~\ref{tab:lablayout} shows what each step costs and
that it works, not speedup. A cold run is the first ever on fresh containers; a
cached run follows it. The last row ran eight jobs at once on containers with
6.0\,GB of memory in total.

\begin{table}[!htbp]
\centering
\caption{End-to-end steps, first run and cached}
\label{tab:lablayout}
\renewcommand{\arraystretch}{1.2}
\rowcolors{2}{}{rowgray}
\begin{tabularx}{\textwidth}{X C{2.4cm} C{2.4cm}}
\toprule
\textbf{Step} & \textbf{Cold} & \textbf{Cached} \\
\midrule
Live join, until connected & 1.8\,s & --- \\
Parallel map & 43.5\,s & 17.0\,s \\
Map over a lambda & 20.5\,s & 4.8\,s \\
GPU job submitted from a GPU-less node & 264.8\,s & 7.6\,s \\
Two-rank GPU ring on one shared GPU & 630.3\,s & 9.5\,s \\
Combined script & 664.6\,s & 26.7\,s \\
Five student-style scripts & 35--126\,s each & 20--70\,s each \\
Concurrent workloads & \multicolumn{2}{c}{412.0\,s, all 8 jobs completed} \\
\bottomrule
\end{tabularx}
\end{table}

Cold times are dominated by building and fetching environments; the GPU steps
download CUDA torch.

\section{Measurements Between Real Machines}
\label{sec:res-recorded}

\textbf{Two laptops in different homes.} Table~\ref{tab:twolaptops} gives runs
between two laptops connected over a direct Tailscale path with a 6--8\,ms round
trip.

\begin{table}[!htbp]
\centering
\caption{Two laptops over the internet (RTX~3060 and RTX~3050)}
\label{tab:twolaptops}
\renewcommand{\arraystretch}{1.2}
\rowcolors{2}{}{rowgray}
\begin{tabularx}{\textwidth}{X L{6.6cm}}
\toprule
\textbf{Run} & \textbf{Result} \\
\midrule
CPU script (pool, joblib, NumPy, pandas) & about 14\,s split, half the items on the peer, against 18\,s on one laptop \\
Whole job sent to the peer & 11.7\,s \\
Small training model, one GPU & 4.2\,s, validation loss 0.0794 \\
Same, \texttt{-{}-ddp~2}, exact synchronisation & 37.1\,s, loss 0.0797 \\
Same, \texttt{-{}-ddp~2}, fast synchronisation & 8.5\,s, loss 0.0797 \\
Same, \texttt{-{}-ddp~2}, automatic choice & 14.6\,s, fast after a short exact warm-up; loss 0.0796 \\
Combined script, cold and warm & 506\,s (487\,s environment build, CUDA torch download); 17.7\,s warm \\
\bottomrule
\end{tabularx}
\end{table}

The mixed CPU script ran faster across the two laptops than on one. The small
training model did not: one GPU finished first. The rings were fixed with
\texttt{-{}-ddp~2}, so the automatic choice could pick only the sync mode, not
release a rank; without \texttt{-{}-ddp~$N$} the lead releases a ring that one
rank alone outruns. Across the internet link, fast synchronisation ran 4.4 times
faster than exact at the same validation loss, and the automatic choice, which
starts exact and switches once, ran 2.5 times faster than exact.

\textbf{Across two NATs.} On the two-site testbed of
Section~\ref{sec:res-setup}, scripts, parallel maps, remote runs and both
training sync modes ran across the sites, and a node whose link was cut mid-run
was replaced. Two machines meet directly when one introducer is reachable;
otherwise a public address or an overlay such as Tailscale is required. That is
the setting of the two-laptop measurements above: neither home offered a
reachable introducer, since one router refuses port mapping and the other
machine sits behind two NATs, so PipedPeer's own links ran over Tailscale. No
relay is involved in either case.

\textbf{A phone.} On a Samsung Galaxy S20~FE (Android~13, Termux), a laptop's
CPU script spread over that laptop, a second laptop over the internet and the
phone gave the phone 8 of 192 parallel items with an unchanged result, after the
phone had built the job's environment from the binary cache in 27\,s. A torch
script training many small models gave the phone a few of them, all converged;
when the laptops could finish first, the phone took none and the result was
unchanged.

\textbf{Further measurements.} Table~\ref{tab:recorded} collects measurements
between a 20-core machine, a 16-core peer and a public cloud host, each taken a
small number of times on the machines and links available.

\begin{table}[!htbp]
\centering
\caption{Further measurements between real machines}
\label{tab:recorded}
\renewcommand{\arraystretch}{1.2}
\rowcolors{2}{}{rowgray}
\small
\begin{tabularx}{\textwidth}{X L{6.2cm}}
\toprule
\textbf{Measurement} & \textbf{Result} \\
\midrule
Parallel map, 20-core machine and 16-core peer over Tailscale, plain Python against PipedPeer & about 6\,s against 5\,s; under half the items on the peer \\
Training on two machines over a home link of about 54\,Mbit/s: one machine / two ranks fp16 / two ranks int8 & 54.9 / 69.4 / 56.1\,s; loss 0.0875 / 0.0875 / 0.0877 \\
Introducer traffic during a cross-network run & 0.2\,MB, against about 30\,MB between the peers \\
Signed closure on a host that requires signatures & imported when the sender's key was trusted; refused when it was removed \\
Daemon memory ceiling, 3\,GiB allocation under a 1\,GiB ceiling & child killed in seconds, out-of-memory recorded \\
\bottomrule
\end{tabularx}
\end{table}

These support modest claims: a parallel map runs faster across two machines
than on one, though short of what the core counts allow; exact training over a
home link does not beat one machine; and the introducer carries a negligible
share of the traffic.

\section{Comparative Analysis}
\label{sec:comparison}

\begin{table}[!htbp]
\centering
\caption{PipedPeer against existing systems. The lower block is where
PipedPeer is partial or the weaker system.}
\label{tab:comparison}
\renewcommand{\arraystretch}{1.2}
\rowcolors{2}{}{rowgray}
\small
\begin{tabularx}{\textwidth}{X C{1.6cm} C{1.2cm} C{1.2cm} C{1.2cm} C{1.35cm} C{1.1cm}}
\toprule
\textbf{Property} & \textbf{Piped\-Peer} & \textbf{Spark} & \textbf{Ray} & \textbf{Dask} & \textbf{BOINC} & \textbf{ssh} \\
\midrule
\multicolumn{7}{l}{\textit{What it does differently}} \\
Runs unmodified Python     & yes & no  & partial & partial & no      & yes \\
No cluster to configure    & yes & no  & no      & no      & no      & partial \\
No head or master node     & yes & no  & no      & no      & no      & yes \\
Reproduces the environment & yes & partial & partial & partial & partial & no \\
Joins NATed machines, no relay$^{\dagger}$ & yes & no & no  & no      & no      & no \\
\midrule
\multicolumn{7}{l}{\textit{Where it is partial or behind}} \\
Survives losing a node$^{*}$ & partial & yes & yes & yes & yes & no \\
Authenticated and encrypted & partial & yes & partial & partial & yes & yes \\
Sandboxes each job          & partial & no  & no      & no      & partial & no \\
Runtimes other than Python  & no & yes & partial & no & yes & yes \\
Proven in production use    & no & yes & yes & yes & yes & yes \\
\bottomrule
\end{tabularx}
\par\smallskip\raggedright\footnotesize $^{\dagger}$When one introducer is
reachable; otherwise a public address or an overlay such as Tailscale is needed.
$^{*}$Any worker; not the submitter of a parallel script or the machine of a
training ring's rank~0.
\end{table}

The upper block lists the criteria on which PipedPeer was designed to differ, so
it wins them by construction; the lower block is the more informative half
(Section~\ref{sec:security}). Spark, Ray and Dask are far more capable at scale,
with mature schedulers and tooling; PipedPeer trades maturity and scale for the
elimination of set-up, code change and network configuration. BOINC
\cite{anderson2004boinc}, closest in spirit, is a platform for specific projects
with their own servers rather than a general executor.
